% Solución explícita y clasificación del universo descendiente.

\section{Explicit Regular Solution of the Descending Cosmological Continuation}
\label{sec:ivv-descendant-explicit}

Considerese the equation efectiva homogenea of Einstein--Cartan for
\(n=3\),
\begin{equation}
 H^2=C\bigl(Aa^{-3}-Ba^{-6}\bigr),
 \qquad A,B,C>0.
 \label{eq:ivv-friedmann-descendant}
\end{equation}
Given \(t_b\in\mathbb R\), define
\begin{equation}
 a(t)^3=\frac BA+\frac{9AC}{4}(t-t_b)^2.
 \label{eq:ivv-scale-factor-descendant}
\end{equation}

\begin{theorem}[Non-Empty Existence and Completeness of the Downward Continuation]
\label{thm:ivv-descendant-existence}
The function \eqref{eq:ivv-scale-factor-descendant} is a global solution of
\eqref{eq:ivv-friedmann-descendant}. It has a unique minimum
\(a_b=(B/A)^{1/3}>0\), satisfies
\[
 H(t_b)=0,
 \qquad
 \dot H(t_b)=\frac{3A^2C}{2B}>0,
\]
and defines on \(\mathbb R\times\mathbb T^3\) a smooth,
causally orientable FLRW metric with bounded curvature invariants and complete
causal geodesics. The regions \(t<t_b\) and \(t>t_b\) realize the contractive
parent branch and the expansive descendant continuation, joined by the
regular hypersurface \(t=t_b\).
\end{theorem}

\begin{proof}
Setting \(y=a^3\), one obtains
\[
 H=\frac{\dot y}{3y}=\frac{3AC(t-t_b)}{2y},
 \qquad
 C(Aa^{-3}-Ba^{-6})
 =C\frac{Ay-B}{y^2}
 =\frac{9A^2C^2(t-t_b)^2}{4y^2}=H^2.
\]
The positivity of \(B/A\) excludes \(a=0\). Both \(H\) and \(\dot H\)
are bounded; the polynomial curvature invariants of the flat FLRW metric
are likewise bounded. Cosmic time ranges over \(\mathbb R\), and the affine
parameter of null geodesics contains \(\int a(t)\,dt\), which diverges at both
ends because \(a(t)\asymp |t|^{2/3}\). The proper temporal length likewise
diverges. There is therefore no finite causal boundary or singularity in the
bounce.
\end{proof}

\begin{theorem}[Obstruction of the classical branched cobordism]
\label{thm:ivv-cobordism-obstruction}
Within the class of smooth, nondegenerate, time-orientable, and globally
hyperbolic Lorentzian manifolds, an evolution between compact Cauchy
surfaces cannot produce branching with a change of topology. Every Cauchy
foliation is diffeomorphic to \(\mathbb R\times\Sigma\); therefore,
the initial and final leaves are diffeomorphic. Branched cobordism is
obstructed within this class.
\end{theorem}

Thus the \(2\) geometric objects are resolved: the minimal continuation
possesses the complete preceding solution, and
the classical branching with topological change is excluded by the obstruction
theorem. A degenerate geometry or a quantum theory of topological
change constitutes another domain and does not modify these results.
